\documentclass[10pt,letterpaper]{article}
\usepackage[utf8]{inputenc}
\usepackage[intlimits]{amsmath}
\usepackage{amsfonts}
\usepackage{amssymb}
\usepackage{ragged2e}
\usepackage[letterpaper, margin=0.5in]{geometry}
\usepackage{graphicx}
\usepackage{cancel}
\usepackage{mathtools}
\usepackage{tabularx}
\usepackage{arydshln}
\usepackage{tensor}
\usepackage{array}
\usepackage{xcolor}
\usepackage[boxed]{algorithm}
\usepackage[noend]{algpseudocode}
\usepackage{listings}
\usepackage{textcomp}
% \usepackage[pdf,tmpdir,singlefile]{graphviz}
\usepackage{mathrsfs}
\usepackage{bbm}
\usepackage{tikz}
\usepackage{tikz-cd}
\usepackage{enumitem}
\usepackage{arydshln}
\usepackage{relsize}
\usepackage{multicol}
\usepackage{scalerel}
\usepackage{upgreek}
\usepackage{ifthen}

\usetikzlibrary{bayesnet}
\setlist{noitemsep}

%%%%%%%%%%%%%%%%%%%%%%%%%%%%%
% Formatting commands
%%%%%%%%%%%%%%%%%%%%%%%%%%%%%
\newcommand{\n}{\hfill\break}
\newcommand{\nn}{\vspace{0.5\baselineskip}\n}
\newcommand{\up}{\vspace{-\baselineskip}}
\newcommand{\hangblock}[2]{\par\noindent\settowidth{\hangindent}{\textbf{#1: }}\textbf{#1: }\nolinebreak#2}
\newcommand{\lemma}[1]{\hangblock{Lemma}{#1}}
\newcommand{\defn}[1]{\hangblock{Defn}{#1}}
\newcommand{\thm}[1]{\hangblock{Thm}{#1}}
\newcommand{\cor}[1]{\hangblock{Cor}{#1}}
\newcommand{\prop}[1]{\hangblock{Prop}{#1}}
\newcommand{\ex}[1]{\hangblock{Ex}{#1}}
\newcommand{\exer}[1]{\hangblock{Exer}{#1}}
\newcommand{\fact}[1]{\hangblock{Fact}{#1}}
\newcommand{\remark}[1]{\hangblock{Remark}{#1}}
\newcommand{\proven}{\;$\square$\n}
\newcommand{\problem}[1]{\par\noindent{\nolinebreak#1}\n}
\newcommand{\problempart}[2]{\par\noindent\indent{}\settowidth{\hangindent}{\textbf{(#1)} \indent{}}\textbf{(#1) }\nolinebreak#2\n}
\newcommand{\ptxt}[1]{\textrm{\textnormal{#1}}}
\newcommand{\inlineeq}[1]{\centerline{$\displaystyle #1$}}
\newcommand{\pageline}{\up\par\noindent\rule{\textwidth}{0.1pt}}

%%%%%%%%%%%%%%%%%%%%%%%%%%%%%
% Math commands
%%%%%%%%%%%%%%%%%%%%%%%%%%%%%
% Set Theory
\newcommand{\card}[1]{\left|#1\right|}
\newcommand{\set}[1]{\left\{#1\right\}}
\newcommand{\setmid}{\;\middle|\;}
\newcommand{\ps}[1]{\mathcal{P}\left(#1\right)}
\newcommand{\pfinite}[1]{\mathcal{P}^{\ptxt{finite}}\left(#1\right)}
\newcommand{\naturals}{\mathbb{N}}
\newcommand{\N}{\naturals}
\newcommand{\integers}{\mathbb{Z}}
\newcommand{\Z}{\integers}
\newcommand{\rationals}{\mathbb{Q}}
\newcommand{\Q}{\rationals}
\newcommand{\reals}{\mathbb{R}}
\newcommand{\R}{\reals}
\newcommand{\complex}{\mathbb{C}}
\newcommand{\C}{\complex}
\newcommand{\halfPlane}{\mathbb{H}}
\let\HSym\H
\let\H\relax
\newcommand{\H}{\halfPlane}
\newcommand{\comp}{^{\complement}}
\DeclareMathOperator{\Hom}{Hom}
\newcommand{\Ind}{\mathbbm{1}}
\newcommand{\cut}{\setminus}
\DeclareMathOperator{\elem}{elem}

% Differentiable Manifolds
\newcommand{\RP}{\mathbb{RP}}
\newcommand{\CP}{\mathbb{RP}}
\newcommand{\osubset}{\overset{\mathclap{\scalebox{0.5}{\ptxt{open}}}}{\subset}}
\newcommand{\osubseteq}{\overset{\mathclap{\scalebox{0.5}{\ptxt{open}}}}{\subseteq}}
\newcommand{\osupset}{\overset{\mathclap{\scalebox{0.5}{\ptxt{open}}}}{\supset}}
\newcommand{\osupseteq}{\overset{\mathclap{\scalebox{0.5}{\ptxt{open}}}}{\supseteq}}
\newcommand{\pdat}[3]{\left.\pd{#1}{#2}\right|_{#3}}
\DeclareMathOperator{\so}{so}
\DeclareMathOperator{\codim}{codim}
\DeclareMathOperator{\Diff}{Diff}
\let\dSym\d
\let\d\relax
\newcommand{\d}{\partial}
\DeclareMathOperator{\gl}{gl}
\DeclareMathOperator{\Vol}{Vol}

% Graph Theory
\let\deg\relax
\DeclareMathOperator{\deg}{deg}
\newcommand{\degp}{\ptxt{deg}^{+}}
\newcommand{\degn}{\ptxt{deg}^{-}}
\newcommand{\precdot}{\mathrel{\ooalign{$\prec$\cr\hidewidth\hbox{$\cdot\mkern0.5mu$}\cr}}}
\newcommand{\succdot}{\mathrel{\ooalign{$\cdot\mkern0.5mu$\cr\hidewidth\hbox{$\succ$}\cr\phantom{$\succ$}}}}
\DeclareMathOperator{\cl}{cl}
\DeclareMathOperator{\affdim}{affdim}

% Probability
\newcommand{\parSymbol}{\P}
\newcommand{\Prob}{\mathbb{P}}
\renewcommand{\P}{\Prob}
\newcommand{\Avg}{\mathbb{E}}
\newcommand{\E}{\Avg}
\DeclareMathOperator{\Var}{Var}
\DeclareMathOperator{\cov}{cov}
\DeclareMathOperator{\Unif}{Unif}
\DeclareMathOperator{\Binom}{Binom}
\newcommand{\CI}{\mathrel{\text{\scalebox{1.07}{$\perp\mkern-10mu\perp$}}}}
\DeclareMathOperator{\Ber}{Ber}
\DeclareMathOperator{\Bin}{Bin}
\DeclareMathOperator{\Geom}{Geom}
\DeclareMathOperator{\Poisson}{Poisson}

% Standard Math
\newcommand{\inv}{^{-1}}
\newcommand{\abs}[1]{\left|#1\right|}
\newcommand{\ceil}[1]{\left\lceil{}#1\right\rceil{}}
\newcommand{\floor}[1]{\left\lfloor{}#1\right\rfloor{}}
\newcommand{\conj}[1]{\overline{#1}}
\newcommand{\of}{\circ}
\newcommand{\tri}{\triangle}
\newcommand{\inj}{\hookrightarrow}
\newcommand{\surj}{\twoheadrightarrow}
\newcommand{\ndiv}{\nmid}
\renewcommand{\epsilon}{\varepsilon}
\newcommand{\divides}{\mid}
\newcommand{\ndivides}{\nmid}
\DeclareMathOperator{\lcm}{lcm}
\DeclareMathOperator{\sgn}{sgn}
\newcommand{\map}[4]{\!\!\!\begin{array}[t]{rcl}#1 & \!\!\!\!\to & \!\!\!\!#2\\ {}#3 & \!\!\!\!\mapsto & \!\!\!\!#4\end{array}}
\newcommand{\bigsum}[2]{\smashoperator[lr]{\sum_{\scalebox{#1}{$#2$}}}}
\DeclareMathOperator{\gcf}{gcf}
\newcommand{\restr}[2]{\left.#1\right|_{#2}}

% Linear Algebra
\newcommand{\Id}{\textrm{\textnormal{Id}}}
\newcommand{\im}{\textrm{\textnormal{im}}}
\newcommand{\norm}[1]{\abs{\abs{#1}}}
\newcommand{\tpose}{^{T}\!}
\newcommand{\iprod}[1]{\left<#1\right>}
\newcommand{\giprod}{\iprod{\;\,,\;}}
\DeclareMathOperator{\tr}{tr}
\DeclareMathOperator{\trace}{tr}
\newcommand{\chgBasMat}[3]{\!\!\tensor*[_{#1}]{\left[#2\right]}{_{#3}}}
\newcommand{\vecBas}[2]{\tensor*[]{\left[#1\right]}{_{#2}}}
\DeclareMathOperator{\GL}{GL}
\DeclareMathOperator{\Mat}{Mat}
\DeclareMathOperator{\vspan}{span}
\DeclareMathOperator{\rank}{rank}
\newcommand{\V}[1]{\vec{#1}}
\DeclareMathOperator{\proj}{proj}
\DeclareMathOperator{\compProj}{comp}
\DeclareMathOperator{\row}{row}
\newcommand{\smallPMatrix}[1]{\paren{\begin{smallmatrix}#1\end{smallmatrix}}}
\newcommand{\smallBMatrix}[1]{\brack{\begin{smallmatrix}#1\end{smallmatrix}}}
\newcommand{\pmat}[1]{\begin{pmatrix}#1\end{pmatrix}}
\newcommand{\bmat}[1]{\begin{bmatrix}#1\end{bmatrix}}
\newcommand{\dual}{^{*}}
\newcommand{\pinv}{^{\dagger}}
\newcommand{\horizontalMatrixLine}{\ptxt{\rotatebox[origin=c]{-90}{$|$}}}
\DeclareMathOperator{\range}{range}
\DeclareMathOperator{\Symm}{Symm}
\DeclareMathOperator{\SU}{SU}

% Multilinear Algebra
\let\Lsym\L
\let\L\relax
\DeclareMathOperator{\L}{\mathscr{L}}
\DeclareMathOperator{\A}{\mathcal{A}}
\DeclareMathOperator{\Alt}{Alt}
\DeclareMathOperator{\Sym}{Sym}
\newcommand{\ot}{\otimes}
\newcommand{\ox}{\otimes}
\DeclareMathOperator{\asc}{asc}
\DeclareMathOperator{\asSet}{set}
\DeclareMathOperator{\sort}{sort}
\DeclareMathOperator{\ringA}{\mathring{A}}
\DeclareMathOperator{\Sh}{Sh}
\DeclareMathOperator{\Bil}{Bil}

% Topology
\newcommand{\closure}[1]{\overline{#1}}
\newcommand{\uball}{\mathcal{U}}
\DeclareMathOperator{\Int}{Int}
\DeclareMathOperator{\Ext}{Ext}
\DeclareMathOperator{\Bd}{Bd}
\DeclareMathOperator{\rInt}{rInt}
\DeclareMathOperator{\ch}{ch}
\DeclareMathOperator{\ah}{ah}
\newcommand{\LargerTau}{\mathlarger{\mathlarger{\mathlarger{\mathlarger{\tau}}}}}
\newcommand{\Tau}{\mathcal{T}}

% Analysis
\DeclareMathOperator{\Graph}{Graph}
\DeclareMathOperator{\epi}{epi}
\DeclareMathOperator{\hypo}{hypo}
\DeclareMathOperator{\supp}{supp}
\newcommand{\lint}[2]{\underset{#1}{\overset{#2}{{\color{black}\underline{{\color{white}\overline{{\color{black}\int}}\color{black}}}}}}}
\newcommand{\uint}[2]{\underset{#1}{\overset{#2}{{\color{white}\underline{{\color{black}\overline{{\color{black}\int}}\color{black}}}}}}}
\newcommand{\alignint}[2]{\underset{#1}{\overset{#2}{{\color{white}\underline{{\color{white}\overline{{\color{black}\int}}\color{black}}}}}}}
\newcommand{\extint}{\ptxt{ext}\int}
\newcommand{\extalignint}[2]{\ptxt{ext}\alignint{#1}{#2}}
\newcommand{\conv}{\ast}
\newcommand{\pd}[2]{\frac{\partial{}#1}{\partial{}#2}}
\newcommand{\del}{\nabla}
\DeclareMathOperator{\grad}{grad}
\DeclareMathOperator{\curl}{curl}
\let\div\relax
\DeclareMathOperator{\div}{div}
\DeclareMathOperator{\vol}{vol}

% Complex Analysis
\let\Re\relax
\DeclareMathOperator{\Re}{Re}
\let\Im\relax
\DeclareMathOperator{\Im}{Im}
\DeclareMathOperator{\Res}{Res}

% Abstract Algebra
\DeclareMathOperator{\ord}{ord}
\newcommand{\generated}[1]{\left<#1\right>}
\newcommand{\cycle}[1]{\smallPMatrix{#1}}
\newcommand{\id}{\ptxt{id}}
\newcommand{\iso}{\cong}
\DeclareMathOperator{\Aut}{Aut}
\DeclareMathOperator{\SL}{SL}
\DeclareMathOperator{\op}{op}
\newcommand{\isom}[4]{\!\!\!\begin{array}[t]{rcl}#1 & \!\!\!\!\overset{\sim}{\to} & \!\!\!\!#2\\ #3 & \!\!\!\!\mapsto & \!\!\!\!#4\end{array}}
\newcommand{\F}{\mathbb{F}}
\newcommand{\acts}{\;\reflectbox{\rotatebox[origin=c]{-90}{$\circlearrowright$}}\;}
\newcommand{\disjunion}{\mathrel{\text{\scalebox{1.07}{$\perp\mkern-10mu\perp$}}}}
\DeclareMathOperator{\SO}{SO}
\DeclareMathOperator{\stab}{stab}
\DeclareMathOperator{\nullity}{nullity}
\DeclareMathOperator{\Perm}{Perm}
\DeclareMathOperator{\nsubgp}{\vartriangleleft}
\DeclareMathOperator{\notnsubgp}{\ntriangleleft}
\newcommand{\presentation}[2]{\left<#1\;\middle|\;#2\right>}
\DeclareMathOperator{\Char}{char}
\DeclareMathOperator{\fchar}{char}
\DeclareMathOperator{\triv}{triv}
\DeclareMathOperator{\reg}{reg}
\DeclareMathOperator{\std}{std}
\DeclareMathOperator{\Func}{Func}
\DeclareMathOperator{\End}{End}

% Convex Optimization
\let\sectionSymbol\S
\let\S\relax
\newcommand{\S}{\mathbb{S}}
\DeclareMathOperator{\dist}{dist}
\DeclareMathOperator{\dom}{dom}
\DeclareMathOperator{\diag}{diag}
\DeclareMathOperator{\ones}{\mathbbm{1}}
\newcommand{\minimizeOver}[3]{\begin{array}{rl}\underset{#1}{\ptxt{minimize}} & #2\\ \ptxt{subject to} & #3\end{array}}
\newcommand{\maximizeOver}[3]{\begin{array}{rl}\underset{#1}{\ptxt{maximize}} & #2\\ \ptxt{subject to} & #3\end{array}}
\newcommand{\minimizationProblem}[2]{\minimizeOver{}{#1}{#2}}
\newcommand{\maximizationProblem}[2]{\maximizeOver{}{#1}{#2}}
\newcommand{\minimizeOverUnconstrained}[2]{\begin{array}{rl}\underset{#1}{\ptxt{minimize}} & #2\end{array}}
\newcommand{\maximizeOverUnconstrained}[2]{\begin{array}{rl}\underset{#1}{\ptxt{maximize}} & #2\end{array}}
\newcommand{\minimizationUnconstrained}[1]{\minimizeOverUnconstrained{}{#1}}
\newcommand{\maximizationUnconstrained}[1]{\maximizeOverUnconstrained{}{#1}}
\DeclareMathOperator{\argmin}{argmin}
\DeclareMathOperator{\argmax}{argmax}

% Proofs
\newcommand{\st}{s.t.}
\newcommand{\unique}{!}
\newcommand{\iffdef}{\overset{\ptxt{def}}{\Leftrightarrow}}
\newcommand{\eqVertical}{\rotatebox[origin=c]{90}{=}}
\newcommand{\mapsfrom}{\mathrel{\reflectbox{\ensuremath{\mapsto}}}}
\newcommand{\mapsdown}{\rotatebox[origin=c]{-90}{$\mapsto$}\mkern2mu}
\newcommand{\mapsup}{\rotatebox[origin=c]{90}{$\mapsto$}\mkern2mu}
\newcommand{\from}{\!\mathrel{\reflectbox{\ensuremath{\to}}}}
\newcommand{\labeledeq}[1]{\overset{\mathclap{\ptxt{#1}}}{=}}
\newcommand{\eqdef}{\labeledeq{def}}

% Brackets
\newcommand{\paren}[1]{\left(#1\right)}
\renewcommand{\brack}[1]{\left[#1\right]}
\renewcommand{\brace}[1]{\left\{#1\right\}}
\newcommand{\ang}[1]{\left<#1\right>}

% Algorithms
\algrenewcommand{\algorithmiccomment}[1]{\hskip 1em \texttt{// #1}}
\algrenewcommand\algorithmicrequire{\textbf{Input:}}
\algrenewcommand\algorithmicensure{\textbf{Output:}}
\newcommand{\algP}{\ptxt{\textbf{P}}}
\newcommand{\algNP}{\ptxt{\textbf{NP}}}
\newcommand{\algNPC}{\ptxt{\textbf{NP-Complete}}}
\newcommand{\algNPH}{\ptxt{\textbf{NP-Hard}}}
\newcommand{\algEXP}{\ptxt{\textbf{EXP}}}
\DeclareMathOperator{\fl}{fl}

%%%%%%%%%%%%%%%%%%%%%%%%%%%%%
% Other commands
%%%%%%%%%%%%%%%%%%%%%%%%%%%%%
\newcommand{\flag}[1]{\textbf{\textcolor{red}{#1}}}
\let\uSym\u
\let\u\relax
\newcommand{\u}[1]{\underline{#1}}
\let\bSym\b
\let\b\relax
\newcommand{\b}[1]{\textbf{#1}}
\let\iSym\i
\let\i\relax
\newcommand{\i}[1]{\textit{#1}}
\let\scSym\sc
\let\sc\relax
\newcommand{\sc}[1]{\textsc{#1}}

%%%%%%%%%%%%%%%%%%%%%%%%%%%%%%%%%%%%%%%
% Make l's curvy in math environments %
%%%%%%%%%%%%%%%%%%%%%%%%%%%%%%%%%%%%%%%
\mathcode`l="8000
\begingroup
\makeatletter
\lccode`\~=`\l
\DeclareMathSymbol{\lsb@l}{\mathalpha}{letters}{`l}
\lowercase{\gdef~{\ifnum\the\mathgroup=\m@ne \ell \else \lsb@l \fi}}%
\endgroup

%%%%%%%%%%%%%%%%%%%%%%%%%
% Fix \vdots and \ddots %
%%%%%%%%%%%%%%%%%%%%%%%%%
\usepackage{letltxmacro}
\LetLtxMacro\orgvdots\vdots
\LetLtxMacro\orgddots\ddots

\makeatletter
\DeclareRobustCommand\vdots{%
	\mathpalette\@vdots{}%
}
\newcommand*{\@vdots}[2]{%
	% #1: math style
	% #2: unused
	\sbox0{$#1\cdotp\cdotp\cdotp\m@th$}%
	\sbox2{$#1.\m@th$}%
	\vbox{%
		\dimen@=\wd0 %
		\advance\dimen@ -3\ht2 %
		\kern.5\dimen@
		% remove side bearings
		\dimen@=\wd2 %
		\advance\dimen@ -\ht2 %
		\dimen2=\wd0 %
		\advance\dimen2 -\dimen@
		\vbox to \dimen2{%
			\offinterlineskip
			\copy2 \vfill\copy2 \vfill\copy2 %
		}%
	}%
}
\DeclareRobustCommand\ddots{%
	\mathinner{%
		\mathpalette\@ddots{}%
		\mkern\thinmuskip
	}%
}
\newcommand*{\@ddots}[2]{%
	% #1: math style
	% #2: unused
	\sbox0{$#1\cdotp\cdotp\cdotp\m@th$}%
	\sbox2{$#1.\m@th$}%
	\vbox{%
		\dimen@=\wd0 %
		\advance\dimen@ -3\ht2 %
		\kern.5\dimen@
		% remove side bearings
		\dimen@=\wd2 %
		\advance\dimen@ -\ht2 %
		\dimen2=\wd0 %
		\advance\dimen2 -\dimen@
		\vbox to \dimen2{%
			\offinterlineskip
			\hbox{$#1\mathpunct{.}\m@th$}%
			\vfill
			\hbox{$#1\mathpunct{\kern\wd2}\mathpunct{.}\m@th$}%
			\vfill
			\hbox{$#1\mathpunct{\kern\wd2}\mathpunct{\kern\wd2}\mathpunct{.}\m@th$}%
		}%
	}%
}
\makeatother

\newcommand{\B}{
	\begin{tikzpicture}
	\filldraw [fill=red, draw=black] (0, 0) rectangle (0.37, 0.45);
	\draw [line width=0.5mm, white ] (0.1,0.08) -- (0.1,0.38);
	\draw[line width=0.5mm, white ] (0.1, 0.35) .. controls (0.2, 0.35) and (0.4, 0.2625) .. (0.1, 0.225);
	\draw[line width=0.5mm, white ] (0.1, 0.225) .. controls (0.2, 0.225) and (0.4, 0.1625) .. (0.1, 0.1);
	\end{tikzpicture}
}

% Allow custom symbols as arrows in tikz-cd
\tikzset{
  symbol/.style={
    draw=none,
    every to/.append style={
      edge node={node [sloped, allow upside down, auto=false]{$#1$}}}
  }
}
% Example: \arrow[r,symbol=\cong]
% https://tex.stackexchange.com/questions/394154/how-to-include-inclusion-subgroup-relationship-in-tikz-cd-diagram

\author{Thomas Cohn}
\title{Math 635 Lecture 5}
\date{1/29/21} % Can also use \today

\begin{document}
\maketitle
\setlength\RaggedRightParindent{\parindent}
\RaggedRight

\par\noindent
Some stuff we can do with Riemannian metrics:
\begin{itemize}
	\item A Riemannian metric on $M$ allows us to define the length of a tangent vector, and the angle between two tangent vectors. For $v,w\in{}T_{p}M$, $\norm{v}=\sqrt{\iprod{v,v}_{p}}$, and $\iprod{v,w}_{p}=\norm{v}\norm{w}\cos\theta$ defines $\theta$, the angle between $v$ and $w$.
	\item A Riemannian metric on $M$ allows us to define the length of a curve in $M$. Say $\gamma:[a,b]\to{}M$. We define its length
	\[
		L(\gamma)=\int_{a}^{b}\norm{\dot\gamma(t)}\,dt=\int_{a}^{b}\sqrt{\iprod{\dot\gamma(t),\dot\gamma(t)}_{\gamma(t)}}\,dt
	\]
	If $\im\gamma\subseteq{}U$, a coordinate patch, with coordinates $(x^{1}(t),\ldots,x^{k}(t))$, then
	\[
		L(\gamma)=\int_{a}^{b}\sqrt{g_{ij}(\gamma(t))\dot{x}^{i}(t)\dot{x}^{j}(t)}\,dt\qquad{}g_{ij}=\iprod{\pd{}{x^{i}},\pd{}{x^{j}}}\in{}C^{\infty}(U)
	\]
\end{itemize}

\lemma{
	$L(\gamma)$ is invariant under re-parameterizations of $\gamma$. Consider $[\alpha,\beta]\overset{\sim}{\underset{t}{\to}}[a,b]\overset{\gamma}{\to}M$. Then
	\[
		\int_{\alpha}^{\beta}\norm{\frac{d\gamma}{ds}(t(s))}\,ds=\int_{a}^{b}\norm{\frac{d\gamma}{dt}(t)}\,dt
	\]
	Proof: The chain rule says $\frac{d\gamma}{ds}=\frac{dt}{ds}\frac{d\gamma}{dt}$. So $\norm{\frac{d\gamma}{ds}}=\norm{\frac{dt}{ds}}\norm{\frac{d\gamma}{dt}}$. Now, use the change of variables formula for integrals.\proven
}

\defn{
	Let $M$ be a connected Riemannian manifold. We define the \u{Riemannian distance function} by, $\forall{}p,q\in{}M$,\n $d(p,q)=\inf\set{L(\gamma)\mid{}\gamma\ptxt{ is a continuous curve, or ``path'', joining $p$ and $q$ that is piecewise smooth}}$.\n
}

\par\noindent
Note: We can replace this definition with just ``smooth'' -- the definitions are equivalent. But that's harder to prove, and this definition will be useful later on.\n

\par\noindent
Note that because $M$ is connected, it's path connected. This means the set of lengths of curves connecting pairs of points is nonempty, so the infimum exists. And because $L(\gamma)\ge{}0$, $d(p,q)\ge{}0$.\n

\thm{
	$d$ is a metric, or distance function, i.e., $\forall{}p,q,r\in{}M$,
	\begin{enumerate}[topsep=0pt, itemsep=0pt, leftmargin=4\parindent, label=(\roman*)]
		\item $d(p,q)\ge{}0$, with $d(p,q)=0\Leftrightarrow{}p=q$
		\item $d(p,q)=d(q,p)$
		\item $d(p,q)+d(q,r)\ge{}d(p,r)$
	\end{enumerate}\up\nn
	Proof: (This is only a partial proof)
	\begin{enumerate}[topsep=0pt, itemsep=0pt, leftmargin=4\parindent, label=(\roman*)]
		\item If $p=q$, take a trivial/constant path. $\dot\gamma\equiv{}0$, so $L(\gamma)=0$. The converse remains to be shown: that $d(p,q)=0$ implies $p=q$. This will be a corollary of the ``Gauss lemma'', which we'll do later on.
		\item We never assumed reparameterizations couldn't reverse the direction of the curve. They can, which directly proves $d(p,q)=d(q,p)$.
		\item Among the paths joining $p$ and $r$ are paths that travel through $q$. Specifically, given any path from $p$ to $q$ and any path from $q$ to $r$, we can concatenate them to get a path from $p$ to $r$.
	\end{enumerate}\up\n
	\proven
}

\par\noindent
Observe: The topology defined by $d$ is the same as the given topology on $M$.\n

\par\noindent
Sometimes, but not always, the infimum is attained, i.e., there exists a minimizing path. In fact, if such a path exists, it's always smooth.\n

\ex{
	Minimizing paths on the sphere are arcs of great circles -- intersections of the sphere with hyperplanes through the origin.\n
}

\ex{
	Minimizing paths don't always exist! Consider $M=\R^{2}\cut\set{0}$. For $p=(43,0)$ and $q=(-43,0)$, $d(p,q)=86$, but there's no path of that length between them (since you can't go through the origin).\n
}

\par\noindent
Refer to Do Carmo, Chapter 1, \sectionSymbol{}2 for more details.\n

\subsection*{Volume Element of an Oriented Riemannian Manifold}

\par\noindent
Reminder: An orientation on an orientable manifold $M$ is determined by a class of top-degree differential forms, $\nu$, with the defining property that $\forall{}p\in{}M$ and any $(v_{1},\ldots,v_{n})$, a positive basis of $T_{p}M$, $\nu_{p}(v_{1},\ldots,v_{n})>0$.\n

\par\noindent
In particular, $\nu$ is nowhere-vanishing. Conversely, a nowhere-vanishing top-degree form can be used to define positive bases, and in turn, an orientation.\n

\defn{
	If $M$ is a orientable Riemannian manifold, its \u{volume form} $\nu$ is defined by the property that $\forall{}p\in{}M$, $\forall(v_{1},\ldots,v_{n})$, a positive, orthonormal basis of $T_{p}M$, one has $\nu_{p}(v_{1},\ldots,v_{n})=1$.\n
}

\par\noindent
Observe: If this condition holds for some positive orthonormal basis, it holds for all positive orthonormal bases. The important calculation is as follows: Fix $p\in{}M$, $(v_{1},\ldots,v_{n})$ a positive orthonormal basis of $T_{p}M$, and $(e_{1},\ldots,e_{n})$ any other ordered basis of $T_{p}M$. Then we can write each $e_{i}=\sum_{l=1}^{n}a_{i}^{l}v_{l}$. For any top-degree form $\nu$, $\nu_{p}(e_{1},\ldots,e_{n})=\det(a_{i}^{l})\nu(v_{1},\ldots,v_{n})$, so if $(e_{1},\ldots,e_{n})$ is also positive and orthonormal, then $\det(a_{i}^{l})=1$.\proven

\subsection*{Computation of the Volume Form in Coordinates}

\par\noindent
Start with $(x_{1},\ldots,x_{n})$, a positive coordinate system with domain $U$. (Recall that this means $\forall{}p\in{}U$, $(\pdat{}{x^{1}}{p},\ldots,\pdat{}{x^{n}}{p})$ is a positive basis of $T_{p}M$.) Apply Gram-Schmidt to each basis (pointwise). We obtain vector fields $v_{1},\ldots,v_{n}$ on $U$ which are orthonormal at each point. And Gram-Schmidt shows the $v$'s are related to the partial derivatives by a smooth matrix, so $\forall{}j$, $v_{j}\in\mathfrak{X}(U)$. And, by possibly permuting the $v_{i}$'s, we can ensure it's a positive basis at every point.\n

\par\noindent
In fact, let's write $\pd{}{x^{i}}=\sum_{l}a_{i}^{l}v_{l}$. Then
\[
	g_{ij}=\iprod{\pdat{}{x^{i}}{p},\pdat{}{x^{j}}{p}}=\sum_{k,l}a_{i}^{l}a_{j}^{k}\underbrace{\iprod{v_{k},v_{l}}}_{=\delta_{k,l}}=\sum_{k}a_{i}^{k}a_{j}^{k}=AA\tpose\quad\ptxt{for}\quad{}A=(a_{i}^{k})
\]
So $\det(g_{ij})=\det(A)^{2}>0$. On the other hand, with our Riemannian volume form $\nu$,
\[
	\nu\paren{\pd{}{x^{1}},\ldots,\pd{}{x^{n}}}=\underbrace{\det(A)}_{\mathclap{=\sqrt{\det(g_{ij})}}}\underbrace{\nu(v_{1},\ldots,v_{n})}_{=1}=\sqrt{\det(g_{ij})}
\]
So in coordinates, $\nu=\sqrt{\det(g_{ij})}\,dx^{1}\wedge\cdots\wedge{}dx^{n}$.\n

\defn{
	The volume of a subset $U$ of a Riemannian manifold is $\Vol(U)=\int_{U}\nu$, where $\nu$ is the Riemannian volume form, if this integral is finite.\n
}

\lemma{
	For any coordinate system $(y^{1},\ldots,y^{n})$ on $U$, positive or not, the Riemannian integral
	\[
		\int_{U}\sqrt{\det\paren{\iprod{\pd{}{y^{i}},\pd{}{y^{j}}}}}\,\underbrace{dy^{1}\cdots{}dy^{n}}_{\mathclap{\ptxt{Riemann integral}}}
	\]
	is equal to $\Vol(U)$.\n
}

\par\noindent
This is true because the change of variables formula for a Riemann integral involves the absolute value of the Jacobian. So in the end, orientation is \i{not} needed to compute volumes of manifolds. In fact, we can even compute volumes of non-orientable manifolds! We generalize by using partitions of unity.

\end{document}